% ===== PRÉAMBULE CANONIQUE — à coller VERBATIM en tête de CHAQUE .tex =====
\documentclass[11pt,a4paper]{article}
\usepackage[utf8]{inputenc}\usepackage[T1]{fontenc}\usepackage{lmodern}
\usepackage[french]{babel}
\usepackage{amsmath,amssymb,mathtools}\usepackage{icomma}
\usepackage[version=4]{mhchem}          % \ce{...} : équations & formules
\usepackage{siunitx}\sisetup{locale=FR} % \qty{}{}, \si{}, \ang{}
\usepackage{chemfig}                    % \chemfig{...} : structures développées
\usepackage{xcolor}\usepackage{colortbl}
\usepackage{tikz}\usepackage{pgfplots}\pgfplotsset{compat=1.18}
\usepackage[most]{tcolorbox}\usepackage{enumitem}\usepackage{array,booktabs}
\usepackage[a4paper,margin=2.2cm]{geometry}
\usetikzlibrary{arrows.meta,calc,angles,quotes,positioning,patterns,backgrounds,shapes.geometric}
\definecolor{primary}{RGB}{222,49,99}\definecolor{primarydark}{RGB}{150,22,55}
\definecolor{defcol}{RGB}{0,114,178}\definecolor{methcol}{RGB}{52,92,148}
\definecolor{excol}{RGB}{0,135,90}\definecolor{attncol}{RGB}{200,40,55}
\tcbset{boxbase/.style={enhanced,breakable,boxrule=0.9pt,arc=2.5pt,
  left=3mm,right=3mm,top=2.3mm,bottom=2.3mm,fonttitle=\bfseries,coltitle=white}}
% --- boîtes TUTORIEL (noms EXACTS, ne pas renommer) ---
\newtcolorbox{cle}[1][]{boxbase,colback=defcol!6,colframe=defcol,
  title={\textsc{À retenir}\ifx\relax#1\relax\relax\else\textmd{ : \itshape #1}\fi}}
\newtcolorbox{methode}[1][]{boxbase,colback=methcol!7,colframe=methcol,sharp corners,
  title={\textsc{Méthode}\ifx\relax#1\relax\relax\else\textmd{ --- \itshape #1}\fi}}
\newtcolorbox{exemplebox}[1][]{boxbase,colback=excol!5,colframe=excol,
  title={\textsc{Exemple}\ifx\relax#1\relax\relax\else\textmd{ : \itshape #1}\fi}}
\newtcolorbox{piege}[1][]{boxbase,colback=attncol!6,colframe=attncol,
  title={\textsc{Piège}\ifx\relax#1\relax\relax\else\textmd{ : \itshape #1}\fi}}
\newtcolorbox{remarque}[1][]{boxbase,colback=black!4,colframe=black!45,
  title={\textsc{Remarque}\ifx\relax#1\relax\relax\else\textmd{ : \itshape #1}\fi}}
% --- boîtes EXERCICE / CORRIGÉ (noms EXACTS, auto-comptées) ---
\newtcolorbox[auto counter]{exo}[1][]{boxbase,colback=primary!4,colframe=primary,
  title={\textsc{Exercice}~\thetcbcounter\ifx\relax#1\relax\relax\else\textmd{ --- \itshape #1}\fi}}
\newtcolorbox[auto counter]{corrige}[1][]{boxbase,colback=excol!4,colframe=excol,
  title={\textsc{Corrigé}~\thetcbcounter\ifx\relax#1\relax\relax\else\textmd{ --- \itshape #1}\fi}}
\newcommand{\R}{\mathbb{R}}\newcommand{\N}{\mathbb{N}}\newcommand{\Z}{\mathbb{Z}}\newcommand{\ds}{\displaystyle}
% (un \usepackage{fancyhdr}+en-tête est possible ; il est ignoré côté web)
% =========================================================================
\begin{document}
\begin{center}\color{primarydark}\large\bfseries Thème 4 — La formule de Nernst \quad$\bullet$\quad Niveau Difficile\end{center}

\begin{exo}[f.é.m.\ complète d'une pile]
Pile $\ce{Zn | Zn^{2+}(0{,}010\,mol/L) || Ag+(0{,}10\,mol/L) | Ag}$ à $25\,^{\circ}$C. Données :
$E_{\mathrm{r}}^{\circ}(\ce{Ag+}/\ce{Ag})=+0{,}80$~V, $E_{\mathrm{r}}^{\circ}(\ce{Zn^{2+}}/\ce{Zn})
=-0{,}76$~V.
\begin{enumerate}[label=\alph*)]
  \item Calcule le potentiel de la cathode (argent, $n=1$).
  \item Calcule le potentiel de l'anode (zinc, $n=2$).
  \item En déduis la f.é.m.\ $\Delta E$.
\end{enumerate}
\end{exo}

\begin{exo}[l'influence du pH sur le permanganate]
Demi-pile \ce{MnO4^- + 8 H+ + 5 e^- <=> Mn^{2+} + 4 H2O} ($E_{\mathrm{r}}^{\circ}=+1{,}51$~V, $n=5$),
avec $[\ce{MnO4^-}]=[\ce{Mn^{2+}}]=0{,}01$~mol/L. On rappelle
$E=E_{\mathrm{r}}^{\circ}+\dfrac{0{,}06}{5}\log\dfrac{[\ce{MnO4^-}]\,[\ce{H+}]^{8}}{[\ce{Mn^{2+}}]}$.
\begin{enumerate}[label=\alph*)]
  \item Calcule $E$ à $\mathrm{pH}=0$ ($[\ce{H+}]=1$~mol/L).
  \item Calcule $E$ à $\mathrm{pH}=3$ ($[\ce{H+}]=10^{-3}$~mol/L).
  \item Que conclure sur le pouvoir oxydant du permanganate quand le milieu devient moins acide ?
\end{enumerate}
\end{exo}

\begin{exo}[remonter à la concentration]
Électrode de cuivre \ce{Cu^{2+} + 2 e^- <=> Cu} ($E_{\mathrm{r}}^{\circ}=+0{,}34$~V, $n=2$). On mesure
$E=+0{,}25$~V. Quelle est la concentration $[\ce{Cu^{2+}}]$ ?
\end{exo}

\begin{exo}[de la f.é.m.\ standard à la constante d'équilibre]
Pour la pile Daniell (couples \ce{Cu^{2+}}/\ce{Cu} et \ce{Zn^{2+}}/\ce{Zn}), la f.é.m.\ standard vaut
$\Delta E^{\circ}=+1{,}10$~V et il s'échange $n=2$ électrons. À l'équilibre $\Delta E=0$, d'où
$\Delta E^{\circ}=\dfrac{0{,}06}{n}\log K$. Estime la constante d'équilibre $K$ de la réaction
$\ce{Zn + Cu^{2+} -> Zn^{2+} + Cu}$ et commente.
\end{exo}

\begin{exo}[une pile de concentration]
On relie deux demi-piles \emph{identiques} \ce{Cu^{2+}}/\ce{Cu} ($E_{\mathrm{r}}^{\circ}=+0{,}34$~V,
$n=2$), l'une avec $[\ce{Cu^{2+}}]=1$~mol/L, l'autre avec $[\ce{Cu^{2+}}]=0{,}01$~mol/L.
\begin{enumerate}[label=\alph*)]
  \item Quelle demi-pile est la cathode ? Justifie.
  \item Calcule la f.é.m.\ $\Delta E$ (alors que $\Delta E^{\circ}=0$).
\end{enumerate}
\end{exo}

\end{document}
